\chapter{A Concrete Naming Certificate for $\TailWindowModulusSynchronizationUp$}
\label{ch:concrete-instances-tail-window-modulus-synchronization-namecert}
\origin{human}

Following Bishop and Bridges on regular Cauchy sequences, the $\TailWindowModulusSynchronizationUp$ packet realizes the regular-Cauchy route of \autoref{sec:real-completeness-regular-cauchy} as a finite synchronization certificate. It sits after the StreamName window surface of \autoref{ch:concrete-instances-streamname-namecert}, the RegSeqRat readback of \autoref{ch:concrete-instances-regseqrat-namecert}, the DyadicRatCore threshold ledger of \autoref{ch:concrete-instances-dyadicratcore-namecert}, the terminal RealUp seal of \autoref{ch:concrete-instances-real-namecert}, and the shared tail threshold discipline of \autoref{ch:concrete-instances-tail-cofinality-budget-namecert}. The packet names the finite synchronization between a tail-window request and a displayed modulus row; it does not choose an infinite subsequence, quotient regular Cauchy names, select a limit, invoke countable choice, or assume ambient Real completeness.

\begin{definition}[Tail-window modulus synchronization carrier]
\label{def:tail-window-modulus-synchronization-carrier}
A $\TailWindowModulusSynchronizationUp$ carrier is a finite $\BHist$ packet
$$
T=(M,B,S,W,D,R,E,H,C,P,N)
$$
with the following displayed rows:
\begin{enumerate}
\item $M$ is the displayed modulus row assigning a unary precision request to a tail threshold.
\item $B$ is the $\TailCofinalityBudgetUp$ row recording that later tail requests refine the same finite threshold route.
\item $S$ is the $\StreamNameUp$ schedule for the regular Cauchy family being inspected.
\item $W$ is the finite tail window selected by $M$ and admitted by $B$.
\item $D$ is the $\DyadicRatCoreUp$ ledger containing the tolerance, threshold comparison, and synchronization scale for $W$.
\item $R$ is the $\RegSeqRatUp$ readback row obtained from the synchronized window and dyadic ledger.
\item $E$ is the $\RealUp$ seal row reached only after the modulus, cofinality, window, dyadic, and regular readback rows have been displayed.
\item $H$ records componentwise $\hsame$ transport for the modulus, cofinality budget, stream schedule, selected window, dyadic ledger, readback, seal, replay, provenance, and naming rows.
\item $C$ records displayed $\Cont$ replay of the route
$$
M,B,S,W,D \longmapsto R \longmapsto E .
$$
\item $P$ is the $\Pkg$ provenance over $\TailWindowModulusSynchronizationUp$, $\TailCofinalityBudgetUp$, $\StreamNameUp$, $\RegSeqRatUp$, $\DyadicRatCoreUp$, $\RealUp$, $\BHist$, $\hsame$, $\Cont$, and $\NameCert$.
\item $N$ is the local $\NameCert_{\TailWindowModulusSynchronizationUp}$ row.
\end{enumerate}
The classifier compares accepted carriers only through $M,B,S,W,D,R,E,H,C,P,N$. It has no coordinate for host sequence equality, quotient real equality, selected limits, countable choice, or ambient completeness of $\RealUp$.
\end{definition}

\begin{theorem}[Tail-window modulus synchronization obligations]
\label{thm:tail-window-modulus-synchronization-obligations}
For every accepted $\TailWindowModulusSynchronizationUp$ carrier $T=(M,B,S,W,D,R,E,H,C,P,N)$, the local $\NameCert_{\TailWindowModulusSynchronizationUp}$ surface consists exactly of carrier admission for the displayed packet; modulus threshold exposure through $M$; cofinal tail-budget exposure through $B$; StreamName window selection through $S,W$; DyadicRatCore synchronization through $D$; RegSeqRat readback through $R$; RealUp seal non-escape through $E$; componentwise transport through $H$; displayed replay through $C$; provenance through $P$; and local naming through $N$.
\end{theorem}
\begin{proof}
Project the carrier. The rows $M,B,S,W,D,R,E,H,C,P,N$ are the only classifier-visible coordinates, so no host sequence equality, quotient real equality, selected limit, choice schedule, or completeness principle can be read from the packet.

Read the synchronization rows in order. The modulus row $M$ chooses the threshold for the displayed precision request, and the tail cofinality budget $B$ records how later requests remain on the same finite route.

The window and readback rows then consume that shared threshold. The StreamName schedule $S$ supplies $W$, the dyadic ledger $D$ records the synchronized tolerance comparison, and the regular readback $R$ is downstream from those rows.

Finally the Real seal and structural rows preserve the route. The row $E$ is available only after $R$, while $H,C,P,N$ transport, replay, source, and name the same finite packet.
\end{proof}

\begin{theorem}[Tail-window modulus synchronization handoff]
\label{thm:tail-window-modulus-synchronization-handoff}
Every Real-completion consumer admitted by an accepted $\TailWindowModulusSynchronizationUp$ packet reads a single finite synchronization route
$$
M,B \longmapsto S,W,D \longmapsto R \longmapsto E .
$$
The route synchronizes the modulus threshold with the tail-window budget before any RegSeqRat readback or RealUp seal is available. It exports no quotient regular Cauchy equality, hidden subsequence selector, selected limit, countable-choice row, or ambient Real completeness principle.
\end{theorem}
\begin{proof}
Start from the regular Cauchy vision site \autoref{sec:real-completeness-regular-cauchy}. A completion-facing read asks for a finite tail window at a displayed precision, so the packet first exposes the modulus row $M$ and the cofinal budget row $B$.

Next the window is synchronized with those rows. The schedule $S$ names the finite observations, $W$ is the selected tail window, and $D$ records the dyadic threshold comparison that keeps the modulus and budget aligned.

The readback and seal are downstream. The row $R$ is obtained from the displayed window and ledger, and the terminal row $E$ is available only after $R$. Since the carrier contains no refused coordinate, the consumer receives exactly the finite synchronization route.
\end{proof}

\falsifiablePrediction{An accepted $\TailWindowModulusSynchronizationUp$ read fails if a Real-completion consumer reaches RegSeqRat readback or a RealUp seal before displaying the modulus row, TailCofinalityBudget row, StreamName window, DyadicRatCore synchronization ledger, hsame transport, Cont replay, Pkg provenance, and local NameCert rows.}

\independenceWitness{tail-cofinality-budget, streamname, regseqrat, dyadicratcore, real: no listed sibling is bijective to $\TailWindowModulusSynchronizationUp$ because this packet binds a modulus threshold and a cofinal tail budget to one synchronized StreamName window, DyadicRatCore ledger, RegSeqRat readback, Real seal, transport, replay, provenance, and local naming surface.}

\closureat{TailWindowModulusSynchronizationUp}{seedStr}

\begin{closurestatus}{\TailWindowModulusSynchronizationUp}
  \constructivestory{A $\TailWindowModulusSynchronizationUp$ packet is generated from a modulus threshold row, a TailCofinalityBudget row, StreamName windows, a DyadicRatCore synchronization ledger, RegSeqRat readback, a RealUp seal, componentwise hsame transport, displayed Cont replay, Pkg provenance, and local naming.}
  \theoryclosure{\seedClosure}
  \scopeclosed{\autoref{def:tail-window-modulus-synchronization-carrier}, \autoref{thm:tail-window-modulus-synchronization-obligations}, and \autoref{thm:tail-window-modulus-synchronization-handoff} seed the finite tail-window modulus synchronization route for regular-Cauchy Real-completion consumers.}
  \formalstatus{\unformalizedV}
  \bridgestatus{none}
  \notclaimed{This chapter does not close quotient regular Cauchy equality, hidden subsequence selection, selected limits, countable choice, ambient Real completeness, Lean verification, or bridge checking.}
  \upgradepath{Obligation closure requires checked carrier admission, modulus threshold exposure, tail cofinality exposure, StreamName window selection, DyadicRatCore synchronization, RegSeqRat readback, Real seal non-escape, transport, replay, provenance, and local naming for BEDC.Derived.TailWindowModulusSynchronizationUp.}
\end{closurestatus}
